\input eplain
\pdfpagewidth=148mm
\pdfpageheight=210mm
\hsize=128mm
\vsize=190mm

\pdfhorigin=0pt
\pdfvorigin=0pt
\hoffset=10mm
\voffset=10mm
\noindent binary tree \leaders\hbox to 10pt {\hss.\hss}  \hfill \xref{node8}
\smallskip
\noindent 4 practical application 2 \leaders\hbox to 10pt {\hss.\hss}  \hfill \xref{node6}
\smallskip
\noindent news \leaders\hbox to 10pt {\hss.\hss}  \hfill \xref{node53}
\smallskip
\noindent 102 binary tree level order traversal \leaders\hbox to 10pt {\hss.\hss}  \hfill \xref{node37}
\smallskip
\noindent 104 max depth binary tree \leaders\hbox to 10pt {\hss.\hss}  \hfill \xref{node9}
\smallskip
\noindent 110 balanced binary tree \leaders\hbox to 10pt {\hss.\hss}  \hfill \xref{node10}
\smallskip
\noindent 133 clone graph \leaders\hbox to 10pt {\hss.\hss}  \hfill \xref{node38}
\smallskip
\noindent 141 linked list cycle \leaders\hbox to 10pt {\hss.\hss}  \hfill \xref{node22}
\smallskip
\noindent 1448 count good nodes \leaders\hbox to 10pt {\hss.\hss}  \hfill \xref{node12}
\smallskip
\noindent 150 evaluate polish notation \leaders\hbox to 10pt {\hss.\hss}  \hfill \xref{node39}
\smallskip
\noindent 155 min stack \leaders\hbox to 10pt {\hss.\hss}  \hfill \xref{node44}
\smallskip
\noindent 15 3sum \leaders\hbox to 10pt {\hss.\hss}  \hfill \xref{node1}
\smallskip
\noindent 167 two sum ii \leaders\hbox to 10pt {\hss.\hss}  \hfill \xref{node15}
\smallskip
\noindent 169 majority element \leaders\hbox to 10pt {\hss.\hss}  \hfill \xref{node27}
\smallskip
\noindent 1 two sum \leaders\hbox to 10pt {\hss.\hss}  \hfill \xref{node0}
\smallskip
\noindent 200 number of islands \leaders\hbox to 10pt {\hss.\hss}  \hfill \xref{node46}
\smallskip
\noindent 206 reverse linked list \leaders\hbox to 10pt {\hss.\hss}  \hfill \xref{node26}
\smallskip
\noindent 207 course schedule \leaders\hbox to 10pt {\hss.\hss}  \hfill \xref{node40}
\smallskip
\noindent 208 implement trie \leaders\hbox to 10pt {\hss.\hss}  \hfill \xref{node41}
\smallskip
\noindent 20 valid parentheses \leaders\hbox to 10pt {\hss.\hss}  \hfill \xref{node16}
\smallskip
\noindent 217 contains duplicate \leaders\hbox to 10pt {\hss.\hss}  \hfill \xref{node31}
\smallskip
\noindent 21 merge two sorted lists \leaders\hbox to 10pt {\hss.\hss}  \hfill \xref{node17}
\smallskip
\noindent 232 implement queue using stacks \leaders\hbox to 10pt {\hss.\hss}  \hfill \xref{node23}
\smallskip
\noindent 235 lowest common ancestor BST \leaders\hbox to 10pt {\hss.\hss}  \hfill \xref{node21}
\smallskip
\noindent 236 lowest common ancestor binary tree \leaders\hbox to 10pt {\hss.\hss}  \hfill \xref{node52}
\smallskip
\noindent 238 product of array except self \leaders\hbox to 10pt {\hss.\hss}  \hfill \xref{node43}
\smallskip
\noindent 242 valid anagram \leaders\hbox to 10pt {\hss.\hss}  \hfill \xref{node18}
\smallskip
\noindent 278 first bad version \leaders\hbox to 10pt {\hss.\hss}  \hfill \xref{node24}
\smallskip
\noindent 322 coin change \leaders\hbox to 10pt {\hss.\hss}  \hfill \xref{node42}
\smallskip
\noindent 33 search in rotated sorted array \leaders\hbox to 10pt {\hss.\hss}  \hfill \xref{node48}
\smallskip
\noindent 383 ransom note \leaders\hbox to 10pt {\hss.\hss}  \hfill \xref{node25}
\smallskip
\noindent 39 combination sum \leaders\hbox to 10pt {\hss.\hss}  \hfill \xref{node49}
\smallskip
\noindent 3 longest substring without repeating characters \leaders\hbox to 10pt {\hss.\hss}  \hfill \xref{node36}
\smallskip
\noindent 409 longest palindrome \leaders\hbox to 10pt {\hss.\hss}  \hfill \xref{node3}
\smallskip
\noindent 425 word squares \leaders\hbox to 10pt {\hss.\hss}  \hfill \xref{node5}
\smallskip
\noindent editorial \leaders\hbox to 10pt {\hss.\hss}  \hfill \xref{node7}
\smallskip
\noindent 437 path sum iii \leaders\hbox to 10pt {\hss.\hss}  \hfill \xref{node13}
\smallskip
\noindent 46 permutations \leaders\hbox to 10pt {\hss.\hss}  \hfill \xref{node50}
\smallskip
\noindent 53 maximum subarray \leaders\hbox to 10pt {\hss.\hss}  \hfill \xref{node32}
\smallskip
\noindent 542 01 matrix \leaders\hbox to 10pt {\hss.\hss}  \hfill \xref{node34}
\smallskip
\noindent 543 diameter of binary tree \leaders\hbox to 10pt {\hss.\hss}  \hfill \xref{node29}
\smallskip
\noindent 560 subarray sum equals k \leaders\hbox to 10pt {\hss.\hss}  \hfill \xref{node14}
\smallskip
\noindent 56 merge intervals \leaders\hbox to 10pt {\hss.\hss}  \hfill \xref{node51}
\smallskip
\noindent 57 insert interval \leaders\hbox to 10pt {\hss.\hss}  \hfill \xref{node33}
\smallskip
\noindent 67 add binary \leaders\hbox to 10pt {\hss.\hss}  \hfill \xref{node28}
\smallskip
\noindent 704 binary search \leaders\hbox to 10pt {\hss.\hss}  \hfill \xref{node19}
\smallskip
\noindent 70 climbing stairs \leaders\hbox to 10pt {\hss.\hss}  \hfill \xref{node2}
\smallskip
\noindent 733 flood fill \leaders\hbox to 10pt {\hss.\hss}  \hfill \xref{node20}
\smallskip
\noindent 79 word search \leaders\hbox to 10pt {\hss.\hss}  \hfill \xref{node4}
\smallskip
\noindent 872 leaf similar trees \leaders\hbox to 10pt {\hss.\hss}  \hfill \xref{node11}
\smallskip
\noindent 876 middle of linked list \leaders\hbox to 10pt {\hss.\hss}  \hfill \xref{node30}
\smallskip
\noindent 973 k closest points origin \leaders\hbox to 10pt {\hss.\hss}  \hfill \xref{node35}
\smallskip
\noindent 98 validate binary search tree \leaders\hbox to 10pt {\hss.\hss}  \hfill \xref{node45}
\smallskip
\noindent 994 rotting oranges \leaders\hbox to 10pt {\hss.\hss}  \hfill \xref{node47}
\smallskip
\vfill \break
\xrdef{node0}
{\medskip\noindent\bf\font\titlefont=cmss17 at 20pt\titlefont 1 two sum}\medskip
{\medskip\noindent\bf\font\titlefont=cmitt12 at 17pt\titlefont 2024-04-17}\medskip
{\smallskip\noindent\bf\font\titlefont=cmr12 at 13pt\titlefont $\bullet$ 21:02: TwoSum}\smallskip
I am a little rusty...
\par

\xrdef{node1}
{\medskip\noindent\bf\font\titlefont=cmss17 at 20pt\titlefont 15 3sum}\medskip
{\medskip\noindent\bf\font\titlefont=cmitt12 at 17pt\titlefont 2024-04-17}\medskip
{\smallskip\noindent\bf\font\titlefont=cmr12 at 13pt\titlefont $\bullet$ 21:51: peeking at ThreeSum}\smallskip
I think I want to try and do all problems related to TwoSum, really get that pattern in my fingers.
\par

{\medskip\noindent\bf\font\titlefont=cmitt12 at 17pt\titlefont 2025-04-19}\medskip
{\smallskip\noindent\bf\font\titlefont=cmr12 at 13pt\titlefont $\bullet$ 16:27: attempting ThreeSum again}\smallskip
The problem hints that if you keep one of the values constant, it turns into a TwoSum problem. So, what you want is to take X, and then find two values Y and Z that add up to -X.
\par

We tried this a few days ago, and it didn't work
\par

{\smallskip\noindent\bf\font\titlefont=cmr12 at 13pt\titlefont $\bullet$ 17:55: Another stab at 3sum}\smallskip
I started reading the editorial, and I solved two sum ii while I was doing this. I think it they want you to sort the list, and then perform it as a two-sum ii problem. When you find the a pair, the ordering is guaranteed.
\par

Having problems with the uniqueness logic.
\par

I think we can take advantage of the sorted list. We only want numbers in ascending order, so make sure the low value starts at the index after the one being checked.
\par

edge case didn't work
\par

Looked at the solution: the twosum problem actually added all three values together, not just one to get a target. It also had logic that skipped over duplicate values. The main loop also broke out early if there were any values greater than zero, it indicated that no numbers following it could sum to zero.
\par

\xrdef{node2}
{\medskip\noindent\bf\font\titlefont=cmss17 at 20pt\titlefont 70 climbing stairs}\medskip
{\medskip\noindent\bf\font\titlefont=cmitt12 at 17pt\titlefont 2024-04-25}\medskip
{\smallskip\noindent\bf\font\titlefont=cmr12 at 13pt\titlefont $\bullet$ 08:05: climbing stairs}\smallskip
I had some time to think abou this while washing dishes
\par

It's definitely dynamic programming. You can definitely build on the previous problems.
\par

when n = 2, you have (1 + 1) or (2).
\par

When n = 3, you have ((1 + 1) + 1), ((2) + 1), (1 + 2)
\par

When n = 4, you have ((1 + 1) + 2)), ((2) + 2), and then ((1 + 1) + 1) + 1), (((2) + 1) + 1), ((1 + 2) + 1).
\par

In other words, s[n] = s[n - 1] + s[n - 2]
\par

In other words, How you get to stair N is one of two ways: your last step is two steps away, or it is one step away. So getting to those steps is already determined from the last iterations.
\par

Okay, it works! I have a weird base case in the beginning that checks for n == 1, I bet I could fix that, if had more time. But I get the gist of it. Moving on.
\par

{\medskip\noindent\bf\font\titlefont=cmitt12 at 17pt\titlefont 2025-04-24}\medskip
{\smallskip\noindent\bf\font\titlefont=cmr12 at 13pt\titlefont $\bullet$ 08:25: climbing stairs (partial)}\smallskip
This one I thought was the dynamic programming problem about stairs, but it's about counting the number of possibilities.
\par

I think it still might be a dynamic programming problem. It has subproblems and substructure.
\par

I'm imagining some kind of loop where you are using the previous answer, then tacking on more answers.
\par

I'm out of time. Will have to think about this later
\par

{\smallskip\noindent\bf\font\titlefont=cmr12 at 13pt\titlefont $\bullet$ 21:51: attempt climbing stairs (aborted)}\smallskip
nvm had a call
\par

\xrdef{node3}
{\medskip\noindent\bf\font\titlefont=cmss17 at 20pt\titlefont 409 longest palindrome}\medskip
{\medskip\noindent\bf\font\titlefont=cmitt12 at 17pt\titlefont 2024-04-25}\medskip
{\smallskip\noindent\bf\font\titlefont=cmr12 at 13pt\titlefont $\bullet$ 08:22: Longest Palindrome (partial)}\smallskip
Adding to my leetcode list.
\par

not enough time to finish this, but will plan.
\par

Palindromes tend to be two-pointer problems. This is problem about finding an optimal value. It asks for the size, not necessarily the palindrome itself.
\par

Longest possible palindrome is N. You can do an initial check on the ends to see if that is possible. Otherwise, you know for sure it can't be N. It might be N - 1.
\par

I was hoping you could whittle down two-pointer style iteratively, but I get the feeling you wouldn't be enumerating combinations properly there.
\par

We could go through each letter and find the next rightmost letter. I think vaguely recall this using a hashmap of sorts. I still had my old solution up which seemed to use a counter. You might be able to be clever then and create a counter, and then use that to check if it's even possible to have a palindrome containing a particular letter (needs to be an even number of letters).
\par

I'm out of time. I'll have to flesh this out more later.
\par

{\medskip\noindent\bf\font\titlefont=cmitt12 at 17pt\titlefont 2025-04-26}\medskip
{\smallskip\noindent\bf\font\titlefont=cmr12 at 13pt\titlefont $\bullet$ 15:50: Longest palindrome (continued)}\smallskip
I think I'm going to try using a hashmap.
\par

Oh! I misread the problem. this is not a palindromatic substring. You can arrange the letters in any order. Oh okay, this is much simpler. All have to do is get a counter, sum all the even values, plus the largest odd value.
\par

Hit an edge case in the submission.
\par

oh I see. You can actually add the odd values, just make it odd -1. If there were any odd values, add 1 at the end.
\par

Okay, I need to take a look at the provided solution for a second. It's a bit more than mine.
\par

Starting backwards, it returns either the whole length of the string if there are no letters with odd counts, or the length of the string, minus an odd character count, plus 1. The loop build up the frequency map as it goes. If the count for that letter at that point in time is odd, increase the odd character count. If even, decrease it.
\par

Actually, approach one (simpler way) is my way. Going to write that one down in my solution.
\par

\xrdef{node4}
{\medskip\noindent\bf\font\titlefont=cmss17 at 20pt\titlefont 79 word search}\medskip
{\medskip\noindent\bf\font\titlefont=cmitt12 at 17pt\titlefont 2025-01-07}\medskip
{\smallskip\noindent\bf\font\titlefont=cmr12 at 13pt\titlefont $\bullet$ 14:56: back to wordsearch}\smallskip
{\smallskip\noindent\bf\font\titlefont=cmr12 at 13pt\titlefont $\bullet$ 15:03: I have reset the code, thinking again about this problem}\smallskip
 
\par

{\smallskip\noindent\bf\font\titlefont=cmr12 at 13pt\titlefont $\bullet$ 15:10: new approach: multiple DFS instead of just one}\smallskip
This is how I'm thinking about it:
\par

iterate through grid, and find top-node candidates
\par

For each top-node, perform DFS to try and find path
\par

return true on first successful path
\par

return false otherwise
\par

{\smallskip\noindent\bf\font\titlefont=cmr12 at 13pt\titlefont $\bullet$ 15:29: traversal logic broken again}\smallskip
This time, on a different edge case. I think my "visited" tracking is broken. I'm marking things as visited whenever they are traversed, but this causes problems if they need to be traversed later. It's the same problem as before.
\par

visited shold only contain nodes in the currently traversed path. once it's no longer being considered in the current position of the current path, it should be returned somehow.
\par

{\smallskip\noindent\bf\font\titlefont=cmr12 at 13pt\titlefont $\bullet$ 15:34: Incredible, the edge cases keep piling up.}\smallskip
This problem is frustrating because you can make small tweaks to solve one edge case, only to run into another. I just moved my "visited" logic a few lines down. It solved one edge case, found another.
\par

{\smallskip\noindent\bf\font\titlefont=cmr12 at 13pt\titlefont $\bullet$ 15:47: Try recursion instead of using a stack}\smallskip
I think it will be easier to implement. I don't think I'm getting the stack logic correct.
\par

{\smallskip\noindent\bf\font\titlefont=cmr12 at 13pt\titlefont $\bullet$ 15:54: My recursive DFS approach works now!}\smallskip
{\smallskip\noindent\bf\font\titlefont=cmr12 at 13pt\titlefont $\bullet$ 15:59: Now, we try for trie.}\smallskip
{\smallskip\noindent\bf\font\titlefont=cmr12 at 13pt\titlefont $\bullet$ 16:14: Trie thoughts so far}\smallskip
Okay, they probably want me to build a trie, then see if the word is in the Trie
\par

But, that can't be right. Building a Trie for all the possibilities seems enormously expensive
\par

Trie is about prefixes, so maybe one gets dynamically built. So, prefixes of the word.
\par

The hint here is that this is used to speed up DFS. So maybe this is a way to keep track of paths found?
\par

I don't know how the trie would be able to keep track of previously visited.
\par

{\smallskip\noindent\bf\font\titlefont=cmr12 at 13pt\titlefont $\bullet$ 16:30: Use Trie as a cache?}\smallskip
If you ignore position, you could absolutely build up a tree that indicates for each letter, what all the directional possibilities are.
\par

{\smallskip\noindent\bf\font\titlefont=cmr12 at 13pt\titlefont $\bullet$ 16:48: I think I might have an idea}\smallskip
The Trie structure here has depth of the length of the word, with each level i corresponding to the letter at index i of word. Begin by populating the root node of the Trie. The children will be the positions of all the nodes with the letter at index 0. Iterate through the children of those nodes, and see which of those positions can connect to the next letter (repeats are okay for now). This happens until there are no more letters or children. From there, attempt to find a match in the Trie, and check to ensure there are no repeats.
\par

{\smallskip\noindent\bf\font\titlefont=cmr12 at 13pt\titlefont $\bullet$ 17:03: And this is where I'll leave it for now. Will implement some other time}\smallskip
Ran out of time. But I'll let the idea sit with me for a bit. The nice thing about this logging system is that it's easier for things to "carry over" into other days.
\par

{\medskip\noindent\bf\font\titlefont=cmitt12 at 17pt\titlefont 2025-01-08}\medskip
{\smallskip\noindent\bf\font\titlefont=cmr12 at 13pt\titlefont $\bullet$ 15:43: Attempts to implement my trie idea}\smallskip
{\smallskip\noindent\bf\font\titlefont=cmr12 at 13pt\titlefont $\bullet$ 16:06: Works in tests, TLE on edge cases}\smallskip
{\smallskip\noindent\bf\font\titlefont=cmr12 at 13pt\titlefont $\bullet$ 16:14: Finally getting around to adding "typing" to graph}\smallskip
I'm going to run this thing locally with the edge case, and you need to import typings explicitely.
\par

{\smallskip\noindent\bf\font\titlefont=cmr12 at 13pt\titlefont $\bullet$ 16:18: The trie is blowing up in size}\smallskip
I wonder if there are duplicates here. I bet there are.
\par

Yup. There are. Going to quickly turn it into a set.
\par

{\smallskip\noindent\bf\font\titlefont=cmr12 at 13pt\titlefont $\bullet$ 16:20: Set unblocks the tree complexity, but that wasn't the bottleneck}\smallskip
{\smallskip\noindent\bf\font\titlefont=cmr12 at 13pt\titlefont $\bullet$ 16:39: Time to stop fighting and look up the answer.}\smallskip
{\smallskip\noindent\bf\font\titlefont=cmr12 at 13pt\titlefont $\bullet$ 16:40: The editorial does not implement trie}\smallskip
{\smallskip\noindent\bf\font\titlefont=cmr12 at 13pt\titlefont $\bullet$ 16:41: This is the wrong problem}\smallskip
The problem is actually word squares. I looked up the wrong problem. This has been a tremendous waste of time.
\par

\xrdef{node5}
{\medskip\noindent\bf\font\titlefont=cmss17 at 20pt\titlefont 425 word squares}\medskip
{\medskip\noindent\bf\font\titlefont=cmitt12 at 17pt\titlefont 2025-01-08}\medskip
{\smallskip\noindent\bf\font\titlefont=cmr12 at 13pt\titlefont $\bullet$ 16:41: This is the wrong problem}\smallskip
The problem is actually word squares. I looked up the wrong problem. This has been a tremendous waste of time.
\par

{\smallskip\noindent\bf\font\titlefont=cmr12 at 13pt\titlefont $\bullet$ 16:49: Getting into word squares}\smallskip
{\smallskip\noindent\bf\font\titlefont=cmr12 at 13pt\titlefont $\bullet$ 16:59: DFS first, then Trie}\smallskip
The card mentions trying to solve with DFS, then to attempt it as a Trie problem. So, we will attempt to solve as a DFS problem.
\par

{\smallskip\noindent\bf\font\titlefont=cmr12 at 13pt\titlefont $\bullet$ 18:01: Tests pass, edge case fixes, more edge cases}\smallskip
{\smallskip\noindent\bf\font\titlefont=cmr12 at 13pt\titlefont $\bullet$ 18:41: I have misread the instructions}\smallskip
The key thing is that the kth row and column read the same string.
\par

I looked at one of the "wrong" answers of mine, and it does work.
\par

I'm going to pick this up tomorrow. Switching gears.
\par

{\medskip\noindent\bf\font\titlefont=cmitt12 at 17pt\titlefont 2025-01-10}\medskip
{\smallskip\noindent\bf\font\titlefont=cmr12 at 13pt\titlefont $\bullet$ 12:36: revisiting problem as DFS}\smallskip
{\smallskip\noindent\bf\font\titlefont=cmr12 at 13pt\titlefont $\bullet$ 12:50: today's intuition}\smallskip
Still doing this as a DFS backtracking problem, where one enumerates a candidate combination, validates, and breaks early for combos that don't work.
\par

A word square can be represented as either a set of rows or a set of columns. A valid word square must work such that row[i]=col[i] for every i from 0 to N, where N is the size of the square (1-4).
\par

Partial candidate solutions are in represented as rows.
\par

To validate the partial candidate, iterate through each row. Ensure that the prefix col[i][0:nrows] matches row[i]. If not return false.
\par

{\smallskip\noindent\bf\font\titlefont=cmr12 at 13pt\titlefont $\bullet$ 13:14: works on solution, but TLE on submission}\smallskip
Attempting to memoize with functools
\par

{\smallskip\noindent\bf\font\titlefont=cmr12 at 13pt\titlefont $\bullet$ 13:27: I believe my DFS is correct, but too slow, time for Trie?}\smallskip
{\smallskip\noindent\bf\font\titlefont=cmr12 at 13pt\titlefont $\bullet$ 13:43: Working out a possible Trie solution}\smallskip
Create Trie Structure from all the words
\par

For each word in word list, append all valid solutions where first row is word.
\par

Use the Trie to enumerate possibilities
\par

{\smallskip\noindent\bf\font\titlefont=cmr12 at 13pt\titlefont $\bullet$ 13:56: You can gradually build word squares from prefixes}\smallskip
For example [ball area lead lady] can be built up incrementally: [b], [ba ar], [bal are lea], [ball area lead lady].
\par

The fact that it's prefixy makes me think I'm in the right direction. Trie traversal is still fuzzy.
\par

{\smallskip\noindent\bf\font\titlefont=cmr12 at 13pt\titlefont $\bullet$ 14:25: Working through by hand Trie traversal}\smallskip
The base example doesn't have multiple possibilities for words down the Trie, so I'm extending these examples to do that. So, in the example word list [area lead wall lady ball], I'm adding "acae" and "lada" in order to get another valid word square  [ball acae lada lead] to go with [ball area lead lady]
\par

{\smallskip\noindent\bf\font\titlefont=cmr12 at 13pt\titlefont $\bullet$ 14:55: I need a break from this}\smallskip
I am now trying to move way from drawing the square visually as I go, and just using the Trie I drew to work out the logic. It's taking some time for things to stick. I'm going to switch to other things now.
\par

{\medskip\noindent\bf\font\titlefont=cmitt12 at 17pt\titlefont 2025-01-13}\medskip
{\smallskip\noindent\bf\font\titlefont=cmr12 at 13pt\titlefont $\bullet$ 17:45: Picking up where I left off.}\smallskip
{\smallskip\noindent\bf\font\titlefont=cmr12 at 13pt\titlefont $\bullet$ 17:51: I am looking up the editorial}\smallskip
\xrdef{node6}
{\medskip\noindent\bf\font\titlefont=cmss17 at 20pt\titlefont 4 practical application 2}\medskip
{\medskip\noindent\bf\font\titlefont=cmitt12 at 17pt\titlefont 2025-01-08}\medskip
{\smallskip\noindent\bf\font\titlefont=cmr12 at 13pt\titlefont $\bullet$ 16:59: DFS first, then Trie}\smallskip
The card mentions trying to solve with DFS, then to attempt it as a Trie problem. So, we will attempt to solve as a DFS problem.
\par

\xrdef{node7}
{\medskip\noindent\bf\font\titlefont=cmss17 at 20pt\titlefont editorial}\medskip
{\medskip\noindent\bf\font\titlefont=cmitt12 at 17pt\titlefont 2025-01-13}\medskip
{\smallskip\noindent\bf\font\titlefont=cmr12 at 13pt\titlefont $\bullet$ 18:29: Back to word square editorial}\smallskip
{\smallskip\noindent\bf\font\titlefont=cmr12 at 13pt\titlefont $\bullet$ 18:54: going into codestudy}\smallskip
\xrdef{node8}
{\medskip\noindent\bf\font\titlefont=cmss17 at 20pt\titlefont binary tree}\medskip
{\medskip\noindent\bf\font\titlefont=cmitt12 at 17pt\titlefont 2025-01-15}\medskip
{\smallskip\noindent\bf\font\titlefont=cmr12 at 13pt\titlefont $\bullet$ 15:35: binary tree outlining}\smallskip
\xrdef{node9}
{\medskip\noindent\bf\font\titlefont=cmss17 at 20pt\titlefont 104 max depth binary tree}\medskip
{\medskip\noindent\bf\font\titlefont=cmitt12 at 17pt\titlefont 2025-01-16}\medskip
{\smallskip\noindent\bf\font\titlefont=cmr12 at 13pt\titlefont $\bullet$ 09:42: leetcode study}\smallskip
As part of my practice here, I want to try outling problems more and chunking them into more components.
\par

I want to get back to trees, and picked this one from the pile.
\par

I decided this problem was too small to be worth doing a code study on, but I did make a components subgraph? We'll see if it's useful at all.
\par

{\medskip\noindent\bf\font\titlefont=cmitt12 at 17pt\titlefont 2025-04-26}\medskip
{\smallskip\noindent\bf\font\titlefont=cmr12 at 13pt\titlefont $\bullet$ 20:32: max depth of binary tree}\smallskip
I think this is just the partial bit to 543 diameter of binary tree. It's recursive DFS.
\par

\xrdef{node10}
{\medskip\noindent\bf\font\titlefont=cmss17 at 20pt\titlefont 110 balanced binary tree}\medskip
{\medskip\noindent\bf\font\titlefont=cmitt12 at 17pt\titlefont 2025-01-16}\medskip
{\smallskip\noindent\bf\font\titlefont=cmr12 at 13pt\titlefont $\bullet$ 10:11: 110}\smallskip
{\smallskip\noindent\bf\font\titlefont=cmr12 at 13pt\titlefont $\bullet$ 10:30: I'm pretty sure I'm misunderstanding the problem}\smallskip
Test case failure. Drawing out the tree.
\par

{\smallskip\noindent\bf\font\titlefont=cmr12 at 13pt\titlefont $\bullet$ 10:31: Nevermind, needed a nonlocal variable}\smallskip
{\smallskip\noindent\bf\font\titlefont=cmr12 at 13pt\titlefont $\bullet$ 10:33: Studying solution, outlining?}\smallskip
{\smallskip\noindent\bf\font\titlefont=cmr12 at 13pt\titlefont $\bullet$ 11:18: That process took way too long}\smallskip
I hope I can get faster, or never look at this problem again.
\par

{\medskip\noindent\bf\font\titlefont=cmitt12 at 17pt\titlefont 2025-04-23}\medskip
{\smallskip\noindent\bf\font\titlefont=cmr12 at 13pt\titlefont $\bullet$ 08:18: balanced binary tree (partial)}\smallskip
Now *this* is a binary tree problem.
\par

I might not have enough time to finish this.
\par

Use DFS to find the max depth of each child node. If they differ by more than one, then they are not balanced.
\par

Think it has to be done recursively. I get the feeling that there could be situations where heights are equal from up high, but not when you go in layer by layer.
\par

To be continued...
\par

{\smallskip\noindent\bf\font\titlefont=cmr12 at 13pt\titlefont $\bullet$ 21:49: Let's try to finish up balanced binary tree}\smallskip
Good lord, I somehow solved it.
\par

Looking the solution anyways.
\par

It's pretty much the same (top-down recursion solution), but maybe a little more concise than mine.
\par

\xrdef{node11}
{\medskip\noindent\bf\font\titlefont=cmss17 at 20pt\titlefont 872 leaf similar trees}\medskip
{\medskip\noindent\bf\font\titlefont=cmitt12 at 17pt\titlefont 2025-01-16}\medskip
{\smallskip\noindent\bf\font\titlefont=cmr12 at 13pt\titlefont $\bullet$ 12:09: Leaf-similar trees}\smallskip
{\smallskip\noindent\bf\font\titlefont=cmr12 at 13pt\titlefont $\bullet$ 12:12: today's intuition}\smallskip
This is a traversal problem. In-order traversal I think? In-order is the one where the leaves will go left-to-right.
\par

The simple way to do this is to produce a function that produces a in-order list of leaves for both trees, then compares the lists. I'm going to do that.
\par

There might be a more clever approach involving comparing both trees at the same time during traversal, but I lack the intuition for that.
\par

{\smallskip\noindent\bf\font\titlefont=cmr12 at 13pt\titlefont $\bullet$ 12:32: My initial approach was correct, editorial is fancy}\smallskip
Doing a codestudy now. Generator/yield is still new to me.
\par

\xrdef{node12}
{\medskip\noindent\bf\font\titlefont=cmss17 at 20pt\titlefont 1448 count good nodes}\medskip
{\medskip\noindent\bf\font\titlefont=cmitt12 at 17pt\titlefont 2025-01-19}\medskip
{\smallskip\noindent\bf\font\titlefont=cmr12 at 13pt\titlefont $\bullet$ 14:48: deeper notetaking}\smallskip
{\smallskip\noindent\bf\font\titlefont=cmr12 at 13pt\titlefont $\bullet$ 14:50: today's intuition}\smallskip
We have a property we are measuring, it is called "good nodes".
\par

We want to count the number of good nodes.
\par

Property of a good node: int path from root to node, there is no node with a value greater than X.
\par

DFS problem, because we are making paths.
\par

Root node is always a good node, so the count will be one if the root is not null, otherwise 0?
\par

DFS can be implemented as a recursive function. I'm thinking we'll need to keep track of the max value as a local variable inside the lexical scope of this function. A counter should be in global scope relative to this function. This counter is always monotonically increasing.
\par

The max value is always the max value seen so far in the recursive path. Recursion adds an implicit stack and keeps track of the max value so far on the path. This could probably be done without a stack, but I'm going to do this recursively first.
\par

I used the range limits to set a min val. Note to self: $-10^4$ is $-10000$, not $-1000$. Heh.
\par

{\smallskip\noindent\bf\font\titlefont=cmr12 at 13pt\titlefont $\bullet$ 15:05: It works. Can I rework it to use an explicit stack now?}\smallskip
My hunch is that we'd use a tupled stack, with the node, and then the max value seen so far. Instead of a built-in function, use a while loop that runs until the stack is empty. Initialize the stack with the root value, assuming it is non-empty.
\par

Instead of making a call to left/right. Push those onto the stack.
\par

Success. Need to be more mindful of control flow in while loop. Don't break or return on empty nodes, just continue.
\par

\xrdef{node13}
{\medskip\noindent\bf\font\titlefont=cmss17 at 20pt\titlefont 437 path sum iii}\medskip
{\medskip\noindent\bf\font\titlefont=cmitt12 at 17pt\titlefont 2025-01-19}\medskip
{\smallskip\noindent\bf\font\titlefont=cmr12 at 13pt\titlefont $\bullet$ 15:13: revisiting}\smallskip
{\smallskip\noindent\bf\font\titlefont=cmr12 at 13pt\titlefont $\bullet$ 15:14: today's intuition}\smallskip
A problem whose return value involves counting things.
\par

The criteria is that there is a downwards path (parent->child) whose value sums to a target value.
\par

Tricky part: the starting point can be from anywhere, not just the root.
\par

{\smallskip\noindent\bf\font\titlefont=cmr12 at 13pt\titlefont $\bullet$ 15:29: idea: work backwards?}\smallskip
The big tricky thing is: how to compensate for paths whose start point can be anywhere? A few things came to mind: keep track of values in queue, sliding window, stack, etc. But these didn't seem to quite fit. Then I thought: why not work backwards. Something something get to the leaf node, then backtrack.
\par

The longest possible path is from a root to a leaf. Because there are positive and negative numbers, it's possible for the sum to go up and down depending on if you remove nodes. Nodes can be removed at both ends.
\par

This isn't feeling quite right. It might be this needs to be more exhaustive.
\par

{\smallskip\noindent\bf\font\titlefont=cmr12 at 13pt\titlefont $\bullet$ 15:35: idea: more brute force. start at top then top + 1, etc}\smallskip
We know the flow goes top down, so this DFS traversal can be done with brute force. Given a node, traverse through its underlying tree, and keep track of the paths that sum to the target sum. At that point, all possible paths with that particular node have been exhausted. Traverse again, this time starting with the children. Continue this until you reach the leaves.
\par

Each traversal should keep track of the summation so far. Inside the DFS function, first check and see if the current summation yields the target. Then, recursively call on children, subtracting the nodes value from the target.
\par

The initial root node (if it exists) will therefore have an initial summation of the root value.
\par

{\smallskip\noindent\bf\font\titlefont=cmr12 at 13pt\titlefont $\bullet$ 15:47: initial attmpt isn't working}\smallskip
{\smallskip\noindent\bf\font\titlefont=cmr12 at 13pt\titlefont $\bullet$ 15:54: got a bit further, broke on edge case}\smallskip
I decided to change my approach slightly. Initially, I was using the accumulative sum, running DFS with the sum + the node, then subtracting it. This next attempt, I set things up so that the DFS would reset the sum, starting with the children. While this works for the two test cases, I ran into edge cases when I submitted.
\par

I'm trying to visualize this traversal as "netting all the combinations". With this new edge case, I'm netting 3 instead of 2.
\par

{\smallskip\noindent\bf\font\titlefont=cmr12 at 13pt\titlefont $\bullet$ 16:05: My traversal is behaving differently than I intended}\smallskip
The current edge case is essentially just a chain of values, and it's actually really helpful because you can see the traversal patterns.
\par

What ends up happening is that the recursion somehow "kicks off" more traversals than I expected.
\par

I have a hunch I shouldn't be calling DFS twice.
\par

{\smallskip\noindent\bf\font\titlefont=cmr12 at 13pt\titlefont $\bullet$ 16:13: Okay that's not working. Trying to visualize the motion I want}\smallskip
It might be that I need to use an explicit DFS stack.
\par

{\smallskip\noindent\bf\font\titlefont=cmr12 at 13pt\titlefont $\bullet$ 16:22: Looking at editorial}\smallskip
{\smallskip\noindent\bf\font\titlefont=cmr12 at 13pt\titlefont $\bullet$ 16:24: Prefix Sum solution. I didn't see it before, and I didn't see it then}\smallskip
It seems may have already racked some hours trying to grok this one. Going to look at this more closely.
\par

{\smallskip\noindent\bf\font\titlefont=cmr12 at 13pt\titlefont $\bullet$ 18:08: codestudy for 437}\smallskip
{\smallskip\noindent\bf\font\titlefont=cmr12 at 13pt\titlefont $\bullet$ 19:17: attempting 437 again}\smallskip
That was fast.
\par

{\smallskip\noindent\bf\font\titlefont=cmr12 at 13pt\titlefont $\bullet$ 19:18: What was my problem before?}\smallskip
Initially, I viewed this as DFS problem where in the traversal, you would traverse from the current node to the end, and then start the process over again for the children. One of the issues was that this traversed nodes too many times. Another issue was that I think it didnt't account for all positive nodes.
\par

The trick was to see this as a prefix sum problem, similar to the problem counting subarrays that sum to a target sum.
\par

My problem was, I had forgotten this prefix sum problem, and couldn't connect the dots.
\par

\xrdef{node14}
{\medskip\noindent\bf\font\titlefont=cmss17 at 20pt\titlefont 560 subarray sum equals k}\medskip
{\medskip\noindent\bf\font\titlefont=cmitt12 at 17pt\titlefont 2025-01-19}\medskip
{\smallskip\noindent\bf\font\titlefont=cmr12 at 13pt\titlefont $\bullet$ 16:54: Doing a codestudy on this problem}\smallskip
\xrdef{node15}
{\medskip\noindent\bf\font\titlefont=cmss17 at 20pt\titlefont 167 two sum ii}\medskip
{\medskip\noindent\bf\font\titlefont=cmitt12 at 17pt\titlefont 2025-04-19}\medskip
{\smallskip\noindent\bf\font\titlefont=cmr12 at 13pt\titlefont $\bullet$ 16:53: studying two-sum ii, in order to grok 3sum}\smallskip
Okay done. The solution was mostly spoiled already for me, but it involved using two-pointer
\par

\xrdef{node16}
{\medskip\noindent\bf\font\titlefont=cmss17 at 20pt\titlefont 20 valid parentheses}\medskip
{\medskip\noindent\bf\font\titlefont=cmitt12 at 17pt\titlefont 2025-04-19}\medskip
{\smallskip\noindent\bf\font\titlefont=cmr12 at 13pt\titlefont $\bullet$ 18:23: Valid Parentheses}\smallskip
I think I remember this one being a stack problem. You push open ones, and then if a closed one appears, it needs to match the last element on the stack.
\par

Issues: got my variables wrong. Didn't check the string length to make sure it was a multiple of 2.
\par

Another edge case: the stack needs to be empty at the end of the function.
\par

\xrdef{node17}
{\medskip\noindent\bf\font\titlefont=cmss17 at 20pt\titlefont 21 merge two sorted lists}\medskip
{\medskip\noindent\bf\font\titlefont=cmitt12 at 17pt\titlefont 2025-04-20}\medskip
{\smallskip\noindent\bf\font\titlefont=cmr12 at 13pt\titlefont $\bullet$ 21:48: merge two sorted lists}\smallskip
linked list problem.
\par

\xrdef{node18}
{\medskip\noindent\bf\font\titlefont=cmss17 at 20pt\titlefont 242 valid anagram}\medskip
{\medskip\noindent\bf\font\titlefont=cmitt12 at 17pt\titlefont 2025-04-22}\medskip
{\smallskip\noindent\bf\font\titlefont=cmr12 at 13pt\titlefont $\bullet$ 07:46: 242 valid anagram}\smallskip
Oops saw my one-liner. This is a hashmap problem, which can be efficiently solved using the Counter type in Python.
\par

\xrdef{node19}
{\medskip\noindent\bf\font\titlefont=cmss17 at 20pt\titlefont 704 binary search}\medskip
{\medskip\noindent\bf\font\titlefont=cmitt12 at 17pt\titlefont 2025-04-22}\medskip
{\smallskip\noindent\bf\font\titlefont=cmr12 at 13pt\titlefont $\bullet$ 07:50: 704 binary search}\smallskip
This problem wants you to implement a divide-and-conquer binary search algorithm. Let's see if I can remember how to set this up.
\par

I have messed it up, damn.
\par

I forgot about adding one to the midpoint index. It solves some of them, but not all of them.
\par

Say, why don't we look at the edge case that's failing, and see if that gives us hints. It's an array of size 1 with the element. Do we add an exception in our binary search? Hopefully, there's a way to adjust the algorithm to naturally compensate.
\par

I keep getting tripped up on the subtleness of binary search.
\par

Okay, I forget how to do binary search. Looking at the answer because I'm at the 15 minute mark.
\par

tweaks: while loop is LT, not LTE, conditional check is LTE, the lo/hi bounds start at 0 and N, not N - 1. Do the loop check when you break out of the loop.
\par

Lots and lots of little things.
\par

\xrdef{node20}
{\medskip\noindent\bf\font\titlefont=cmss17 at 20pt\titlefont 733 flood fill}\medskip
{\medskip\noindent\bf\font\titlefont=cmitt12 at 17pt\titlefont 2025-04-22}\medskip
{\smallskip\noindent\bf\font\titlefont=cmr12 at 13pt\titlefont $\bullet$ 08:14: Flood fill (planned, but didn't finish)}\smallskip
It's traversal? Use a queue (or stack?) to traverse directions. Use a set to keep track of where you've been (coordinates can work in a set because tuples are hashable). Get the color of the starting pixel, push the starting pixel onto the queue. While the queue has things, process.
\par

Inside the process: check the 4 directions of the current position. If the pixel hasn't been visited before and is the color of the start pixel color, switch it to the new color, mark it as visited in the set, and append it to the queue.
\par

Make sure the mark the start pixel as visited.
\par

Queue I think needs to be used because it's BFS. I think the order of processing matters? Flood fill should move outwards.
\par

{\smallskip\noindent\bf\font\titlefont=cmr12 at 13pt\titlefont $\bullet$ 19:48: flood fill continued}\smallskip
Well, something is wrong with my implementation
\par

return image was in the wrong place, now the base cases pass
\par

It works!
\par

\xrdef{node21}
{\medskip\noindent\bf\font\titlefont=cmss17 at 20pt\titlefont 235 lowest common ancestor BST}\medskip
{\medskip\noindent\bf\font\titlefont=cmitt12 at 17pt\titlefont 2025-04-23}\medskip
{\smallskip\noindent\bf\font\titlefont=cmr12 at 13pt\titlefont $\bullet$ 07:47: LCA}\smallskip
Binary tree problem. DFS.
\par

"lowest" means closest to leaves.
\par

What are the possibilities?
\par

One is that one of the nodes is an ancestor of the other node, then you return the node that is highest.
\par

Otherwise, the nodes have a common ancestor.
\par

I guess an approach would be: find the first node in the tree, then search for the other node, backtracking if needed.
\par

15 minutes passed. Strugglig formulating an answer. I don't think my thinking is right. Giving up early.
\par

Oh, it's a binary *search* tree. Glossed over that bit.
\par

Wrote up algorithm solution. Typing out python solution now to get it in my fingers.
\par

\xrdef{node22}
{\medskip\noindent\bf\font\titlefont=cmss17 at 20pt\titlefont 141 linked list cycle}\medskip
{\medskip\noindent\bf\font\titlefont=cmitt12 at 17pt\titlefont 2025-04-24}\medskip
{\smallskip\noindent\bf\font\titlefont=cmr12 at 13pt\titlefont $\bullet$ 07:43: Linked list cycle}\smallskip
Linked list problem. Floyd's algorithm. Two moving pointers, one fast, one slow. The trick will be making sure I can get the pointer stuff set up right.
\par

It passes, but I don't like how my while loop has to check in three places (slow, fast, fast.next). Checking my work.
\par

Comparing solution to my answer.
\par

"fast" is head.next, not head
\par

while loop is fast != slow
\par

fast/fast.next checks are done in the while loop
\par

It looks like it's mostly equivalent, just the comparisons are done in a different place?
\par

Their function returns True, mine returns false.
\par

\xrdef{node23}
{\medskip\noindent\bf\font\titlefont=cmss17 at 20pt\titlefont 232 implement queue using stacks}\medskip
{\medskip\noindent\bf\font\titlefont=cmitt12 at 17pt\titlefont 2025-04-24}\medskip
{\smallskip\noindent\bf\font\titlefont=cmr12 at 13pt\titlefont $\bullet$ 07:56: 232 implement queue using stacks}\smallskip
Stack: LIFO. Queue: FIFO. There's no requirements for efficiency
\par

Important: instructions say to only use two stacks
\par

There was something about popping all items then pushing them back on again.
\par

Push is going to be the inefficient operation. We push a new item to stack 2, then drain s1 to s2, then populate s1 again.
\par

Actually,more order was wrong. Drain s1 to s2 to reverse the order, push to s2, then drain s2 to s1 to get the order back.
\par

\xrdef{node24}
{\medskip\noindent\bf\font\titlefont=cmss17 at 20pt\titlefont 278 first bad version}\medskip
{\medskip\noindent\bf\font\titlefont=cmitt12 at 17pt\titlefont 2025-04-24}\medskip
{\smallskip\noindent\bf\font\titlefont=cmr12 at 13pt\titlefont $\bullet$ 08:09: First bad version}\smallskip
This is literally git bisect isn't it?
\par

Divide and conquer. We want to make sure we get the first *bad* version, not the first *good* version, so there migth be some n+1 things happening.
\par

Turned out to be no +1. hi = mid, low = mid + 1 in the checks.
\par

Sorry, this is called "binary search", "divide and conquer" is just slang I think.
\par

\xrdef{node25}
{\medskip\noindent\bf\font\titlefont=cmss17 at 20pt\titlefont 383 ransom note}\medskip
{\medskip\noindent\bf\font\titlefont=cmitt12 at 17pt\titlefont 2025-04-24}\medskip
{\smallskip\noindent\bf\font\titlefont=cmr12 at 13pt\titlefont $\bullet$ 08:16: Ransom Note}\smallskip
I do not remember what is involved in this one at all
\par

Oh, it's not in my leetcode files either
\par

Looks like a hashmap problem. You want to get a frequency map for both the magazine and the ransom letter. For each letter in the random letter, you want to make sure that the magazine has enough. If it doesn't exist in magazine or there's not enough, return false.
\par

Yup, that's how you do it.
\par

\xrdef{node26}
{\medskip\noindent\bf\font\titlefont=cmss17 at 20pt\titlefont 206 reverse linked list}\medskip
{\medskip\noindent\bf\font\titlefont=cmitt12 at 17pt\titlefont 2025-04-26}\medskip
{\smallskip\noindent\bf\font\titlefont=cmr12 at 13pt\titlefont $\bullet$ 16:13: reverse linked list}\smallskip
This one I think is about getting the flow right.
\par

For any given node, their "next" pointer needs to point to the "previous" node. The old "next" pointer entry becomes the current node, and the old current node becomes the previous.
\par

How to set this up? How to return properly?
\par

I'm still trying to make prev/node setup work, but it's not terminating properly.
\par

Okay, let's start again. What happens to make this terminate, and what do we expect?
\par

I forget how to reverse a linked list. Going to chew on this for about 5 more minutes, then I look up the answer.
\par

Wow, okay I did it. I blew all my work up, and started building it up bit by bit. First, just iterate through the linked list. We know we need to return the last value, so keep track of the previous value. We also know that the "next" pointer of the current node needs to point to the previous node. There's a temporary value needed to keep track of the next node. There done.
\par

\xrdef{node27}
{\medskip\noindent\bf\font\titlefont=cmss17 at 20pt\titlefont 169 majority element}\medskip
{\medskip\noindent\bf\font\titlefont=cmitt12 at 17pt\titlefont 2025-04-26}\medskip
{\smallskip\noindent\bf\font\titlefont=cmr12 at 13pt\titlefont $\bullet$ 16:41: majority element}\smallskip
I mean, naively. This seems to be one you can throw a frequency map at, right?
\par

Okay, I solved it. But what does the solution want you to do? It feels like it was asking for something else.
\par

nope, that's the answer. use a hashmap.
\par

The solution was very nice and pythonic: === class Solution:     def majorityElement(self, nums: List[int]) -> int:         counts = Counter(nums)         return max(counts.keys(), key=counts.get) ===
\par

\xrdef{node28}
{\medskip\noindent\bf\font\titlefont=cmss17 at 20pt\titlefont 67 add binary}\medskip
{\medskip\noindent\bf\font\titlefont=cmitt12 at 17pt\titlefont 2025-04-26}\medskip
{\smallskip\noindent\bf\font\titlefont=cmr12 at 13pt\titlefont $\bullet$ 16:48: add binary}\smallskip
This one is silly to do in python because they have built-ins for this. Let me cheat first.
\par

Here's the one-liner === class Solution:     def addBinary(self, a: str, b: str) -> str:         return bin(int(a, 2) + int(b,2))[2:] === What they probably want to do is add the bits one by one. Converting each one into an array turns them into a stack, where the last value for both are the ones place. You could pop values, add, and push the results. Use a carry, and tack on a one at the end if there's still a carry. Reverse and convert to to string.
\par

I can never remember how to convert a list to a string in python
\par

My logic is still wrong. The base cases pass though.
\par

My carry logic was wrong. The overflow can be either 2 (1 + 1) or 3 (1 + 1 + previous carry). When it's three, the current bit is set to be 1, otherwise 2.
\par

\xrdef{node29}
{\medskip\noindent\bf\font\titlefont=cmss17 at 20pt\titlefont 543 diameter of binary tree}\medskip
{\medskip\noindent\bf\font\titlefont=cmitt12 at 17pt\titlefont 2025-04-26}\medskip
{\smallskip\noindent\bf\font\titlefont=cmr12 at 13pt\titlefont $\bullet$ 19:56: Diameter of a Binary Tree}\smallskip
Binary tree. At first it sounded like BFS, but I think it's actually DFS.
\par

It feels like you should traverse to the leftmost node of a tree, then the rightmost node of a tree.
\par

This feels recursive, like you can get keep getting the max diameter of smaller and smaller trees, and building up from that.
\par

You might not be able to re-use results. They did say it didn't have to pass through root.
\par

Solved! I kind of walked myself through the formula. The diameter is the maximum value of something, so that means storing some variable that gets updated through the traversal. Traverse starting at root, return that variable. A maximum diameter means adding the max height of the left/right child node. You can find the max height through recursion.
\par

\xrdef{node30}
{\medskip\noindent\bf\font\titlefont=cmss17 at 20pt\titlefont 876 middle of linked list}\medskip
{\medskip\noindent\bf\font\titlefont=cmitt12 at 17pt\titlefont 2025-04-26}\medskip
{\smallskip\noindent\bf\font\titlefont=cmr12 at 13pt\titlefont $\bullet$ 20:26: middle of a linked list}\smallskip
Linked list problem. I believe this is best solved using fast/slow pointers. Fast pointer moves at 2x speed. When the fast reaches the end return the slow pointer.
\par

I had a little bit of trouble figuring where to start the fast pointer. fast/slow both start at head.
\par

\xrdef{node31}
{\medskip\noindent\bf\font\titlefont=cmss17 at 20pt\titlefont 217 contains duplicate}\medskip
{\medskip\noindent\bf\font\titlefont=cmitt12 at 17pt\titlefont 2025-04-26}\medskip
{\smallskip\noindent\bf\font\titlefont=cmr12 at 13pt\titlefont $\bullet$ 20:37: contains duplicate}\smallskip
Shoot, I saw the answer. It was a one-liner. You turn the array into a set, and compare lengths with the original array.
\par

\xrdef{node32}
{\medskip\noindent\bf\font\titlefont=cmss17 at 20pt\titlefont 53 maximum subarray}\medskip
{\medskip\noindent\bf\font\titlefont=cmitt12 at 17pt\titlefont 2025-04-26}\medskip
{\smallskip\noindent\bf\font\titlefont=cmr12 at 13pt\titlefont $\bullet$ 20:40: maximum subarray}\smallskip
Which one is the subarray? is that the contiguous one or the other one?
\par

Oh good contiguous. That one is a little bit easier I think.
\par

I'm getting a sliding window feel from this? I don't think it's 2-pointer. It might be greedy.
\par

I'm trying greedy. we go left to right and have a running total starting with the first value, and a value indicating the leftmost value in the sum. Every loop we determine what the local optimum is based on the current value. If we add the value and the sum is greater, we include it. If it's less, we start a new sum with the current value. If subtracting the leftmost value makes the sum greater, increment the leftmost value pointer.
\par

Not feeling great about this one. I tried the greedy approach and it didn't work.
\par

oh, the topics are array, divide and conquer, and dynamic programming. The follow-up is divide and conquer.
\par

Okay if it's dynamic programming, we know how to set this up.
\par

It's hot. I think I'll return to this one later.
\par

Generally you're trying to find some max start/end range. When the start/end are the same, there is only one value.
\par

Tried forcing this into a 2d dynamic programming problem. Wrong turn. It's Kadane's algorithm. I will revisit this tomorrow.
\par

{\medskip\noindent\bf\font\titlefont=cmitt12 at 17pt\titlefont 2025-04-28}\medskip
{\smallskip\noindent\bf\font\titlefont=cmr12 at 13pt\titlefont $\bullet$ 07:55: study answer to maximum subarray}\smallskip
initialize 2 variables 2 integer variables currentSubarray and maxSubarray, set them equal to first value in subarray
\par

iterate through the array starting with 2nd element. Add elemen to currentSubarray. Choose the greater value to update currentSubarray: currentSubarray, or the current value.
\par

Update max value.
\par

My greedy intuition was close, but too complicated.
\par

\xrdef{node33}
{\medskip\noindent\bf\font\titlefont=cmss17 at 20pt\titlefont 57 insert interval}\medskip
{\medskip\noindent\bf\font\titlefont=cmitt12 at 17pt\titlefont 2025-04-28}\medskip
{\smallskip\noindent\bf\font\titlefont=cmr12 at 13pt\titlefont $\bullet$ 08:06: Insert Interval}\smallskip
These interval problems always seem so hyper-specific, which makes them feel intimidating to me.
\par

Intervals are sorted and non-overlapping, and you want to insert an interval somewhere in there, and make sure it returns intervals that are still sorted and non-overlapping.
\par

The instructions say we can make a copy, which makes some of the logistics easier I think.
\par

I'm going to start by just setting up the scaffolding that just makes a copy of the existing intervals without copying. I'm going to try to add the insertion logic on top of that.
\par

An insertion interval is about bounds checking for start and end times. Working backwards, we know we have passed the merge when the start time of the current interval is passed the end time of the inserted interval.
\par

What are the possibilities for insertion? One is it inserts without merging, else it merges.
\par

For merging, the inserted interval can overlap before or after an interval. It can be absorbed by an interval. It can also end up merging many intervals.
\par

When is a merge first triggered? A merge is triggered when there is an initial intersection. The current interval overlaps with the inserted interval. The inserted interval starts before and ends before, starts after and ends before, or starts after and ends after.
\par

If we've already started a merge, we need to determine if the current interval needs to be merged in as well, or if we are done merging.
\par

Not sure how we are supposed to handle cases where it doesn't merge.
\par

to be continued...
\par

{\smallskip\noindent\bf\font\titlefont=cmr12 at 13pt\titlefont $\bullet$ 21:28: ...continued}\smallskip
Managed to get this working. I had to blow up my morning thoughts, whip out a piece of paper and pen, and draw a picture.
\par

The key shift in my thinking for this was: an inserted interval is always going to be a merged interval. Another thing I did to help me was to destructure these intervals so I could see them labelled as start/end.
\par

Somehow, I was able to see that the result of merging two intervals together was always going to be the smallest of the start times plus the largest of the end times. And you'd keep doing that with intervals until the start time of the current interval appeared after the merged interval.
\par

A merged interval gets merged when the start time of the interval after it is past the end time of the merged interval, or there are no more intervals. Intervals after the merged interval can be merged normally.
\par

It the current interval's end time is less than the start time of the merged interval, the merged interval has not yet started, and you can just intervals.
\par

\xrdef{node34}
{\medskip\noindent\bf\font\titlefont=cmss17 at 20pt\titlefont 542 01 matrix}\medskip
{\medskip\noindent\bf\font\titlefont=cmitt12 at 17pt\titlefont 2025-04-29}\medskip
{\smallskip\noindent\bf\font\titlefont=cmr12 at 13pt\titlefont $\bullet$ 07:48: 01 Matrix}\smallskip
Getting BFS vibes.
\par

The problem is asking for closest distances from 0 tile. If you are a 0 tile, the distance is 0. Anything sharing an edge with a 0 tile that is not a 0 tile will have a distance of 1, and then 2, 3, etc.
\par

We can initialize the output array by iterating through the entire matrix, and finding all the zero tiles, marking them as having distance 0. These can get pushed onto a queue with their location, and their distance (0).
\par

Processing happens by popping items off the queue and traversing in the 4 directions. If a node hasn't been visited, it's distance is marked based on the current distance (+ 1), and then that's pushed onto the queue.
\par

Done. Pretty formulaic BFS. It's definitely a pattern I've learned doing this leetgrind.
\par

\xrdef{node35}
{\medskip\noindent\bf\font\titlefont=cmss17 at 20pt\titlefont 973 k closest points origin}\medskip
{\medskip\noindent\bf\font\titlefont=cmitt12 at 17pt\titlefont 2025-04-29}\medskip
{\smallskip\noindent\bf\font\titlefont=cmr12 at 13pt\titlefont $\bullet$ 08:04: K closest points to origin}\smallskip
This is a problem I think you can solve with a heap. Basically, you push points onto the heap, and take the first K values as the answer. You can also keep a fixed sized heap of K items, and pop the largest value.
\par

python's heapq is a minheap, so I think you'd need to negate distances to keep the smallest ones on the queue the longest.
\par

I am running into a weird edge case with one of the tests.
\par

I was popping before pushing, and should have been pushing before popping.
\par

\xrdef{node36}
{\medskip\noindent\bf\font\titlefont=cmss17 at 20pt\titlefont 3 longest substring without repeating characters}\medskip
{\medskip\noindent\bf\font\titlefont=cmitt12 at 17pt\titlefont 2025-04-29}\medskip
{\smallskip\noindent\bf\font\titlefont=cmr12 at 13pt\titlefont $\bullet$ 08:22: Longest substring without repeating characters (partial attempt)}\smallskip
You're not finding *which* substring, just the max, so this is an optimization problem. This means that keeping track of a max value and returning that.
\par

Feels kinda sliding window, also hashmap. You could probably sweep through left to right, use the map to keep track of indices where the letter was previously found.
\par

Okay, this isn't quite right. I think my hashmap thinking is correct, but I have the logic wrong for updating the start index.
\par

I have to go.
\par

{\smallskip\noindent\bf\font\titlefont=cmr12 at 13pt\titlefont $\bullet$ 20:43: another stab at it}\smallskip
Another idea: use the hashmap. If there is an entry already there, subtract the difference. Otherwise, start from 0.
\par

not even close. looking at the answer.
\par

Answer is "HashMap with Sliding window" (they erroneously call it a HashSet, pretty sure that's wrong)
\par

Set up an empty counter/hashmap and left/right pointers equal to 0.
\par

Loop while right pointer is less than string length.
\par

1: get the character at the right pointer. 2: update the count for that character. 3: while the count of that right character is greater than 1, get the value of the left pointer, decrease that from the counter, and increase the left pointer. 4: obtain the max value of right - left + 1. 5: increment right pointer.
\par

\xrdef{node37}
{\medskip\noindent\bf\font\titlefont=cmss17 at 20pt\titlefont 102 binary tree level order traversal}\medskip
{\medskip\noindent\bf\font\titlefont=cmitt12 at 17pt\titlefont 2025-04-30}\medskip
{\smallskip\noindent\bf\font\titlefont=cmr12 at 13pt\titlefont $\bullet$ 07:46: binary tree level order traversal}\smallskip
This is a BFS tree traversal problem, so use a queue.
\par

I've done double queues before, which can be a little bit tricky to set up. A single queue keeping track of level should suffice.
\par

Basically keep pushing things on the queue, children are always +1. append things to a list. when the next items level is greater than the current level, append the list to the output, and start a new list.
\par

I got it working. I don't love that I have to do a last check at the end of the loop. I'm sure that could be improved, but it's good enough.
\par

The solution just uses the level number directly in an array.
\par

\xrdef{node38}
{\medskip\noindent\bf\font\titlefont=cmss17 at 20pt\titlefont 133 clone graph}\medskip
{\medskip\noindent\bf\font\titlefont=cmitt12 at 17pt\titlefont 2025-04-30}\medskip
{\smallskip\noindent\bf\font\titlefont=cmr12 at 13pt\titlefont $\bullet$ 08:05: clone graph}\smallskip
This is some kind of traversal, probably BFS.
\par

Starting with the first node, go through all the children. Use a lookup table using node values, which are unique. If the node has already been cloned, use that cloned value from the list, otherwise clone the reference, add it to the list, and update the node.
\par

Right, I am scewing up the traversal here. I am trying to get the BFS to both initialize a node, and possibly the neighbor nodes, in the nodelist I am making. But both of them can already have been created.
\par

The solution uses recursion, I use a queue, not sure why mine isn't working.
\par

I'm out of time. Will revisit later.
\par

{\medskip\noindent\bf\font\titlefont=cmitt12 at 17pt\titlefont 2025-05-08}\medskip
{\smallskip\noindent\bf\font\titlefont=cmr12 at 13pt\titlefont $\bullet$ 07:59: clone graph, continued}\smallskip
I did not finish this one
\par

It looks like I forgot to copy over the values themselves, which may be why I was getting the error.
\par

I was appending the wrong thing to the queue. I should have been the neighbor, not the newneighbor.
\par

\xrdef{node39}
{\medskip\noindent\bf\font\titlefont=cmss17 at 20pt\titlefont 150 evaluate polish notation}\medskip
{\medskip\noindent\bf\font\titlefont=cmitt12 at 17pt\titlefont 2025-05-08}\medskip
{\smallskip\noindent\bf\font\titlefont=cmr12 at 13pt\titlefont $\bullet$ 08:15: polish notation}\smallskip
I remember something weird was tripping me up with this one. Might have been the thing that converts strings to numbers? Or possibly the division?
\par

This is basically a stack problem.
\par

I got tripped up again. divide towards zero means int(a/b), not a // b.
\par

\xrdef{node40}
{\medskip\noindent\bf\font\titlefont=cmss17 at 20pt\titlefont 207 course schedule}\medskip
{\medskip\noindent\bf\font\titlefont=cmitt12 at 17pt\titlefont 2025-05-09}\medskip
{\smallskip\noindent\bf\font\titlefont=cmr12 at 13pt\titlefont $\bullet$ 07:59: course schedule}\smallskip
I remember this problem. It's a DAG problem.
\par

This is sort of a traversal problem? You can take all the prereqs, and put them in an adjacency list: a lookup table for each node, and all their rerequisites.
\par

after constructing the adjacency list, traverse BFS using a queue. Begin by adding all the courses that do not have prereqs onto the queue. Pop items off the queue, and remove that course dependency from the adjacency list. Add any courses that "free up".
\par

I'm realizing I could construct my adjacency list a little bit differently. Originally, I wanted to have the list show the node and their prereqs, but the issue with that is that you'd need to do many updates every time a node is removed. It would be easier to show for each node, which courses are dependent on it.
\par

With this new list, you'd need to work backwards. You'd add all the courses that are not dependencies for any other courses, which would be the last courses.... wait this isn't going to work.
\par

Okay, back to the original plan I think. Maybe maintain two lists?
\par

or just have a counter that keeps track of the number of
\par

we are out of time. I can't visualize the traversal correctly.
\par

{\smallskip\noindent\bf\font\titlefont=cmr12 at 13pt\titlefont $\bullet$ 08:43: revisiting}\smallskip
I was being do fancy, just pop dependency from set, do a full sweep
\par

oh shoot, this is the one where you can add phantom edges isn't it.
\par

Got it! Will clean up code and type up solution later
\par

\xrdef{node41}
{\medskip\noindent\bf\font\titlefont=cmss17 at 20pt\titlefont 208 implement trie}\medskip
{\medskip\noindent\bf\font\titlefont=cmitt12 at 17pt\titlefont 2025-05-14}\medskip
{\smallskip\noindent\bf\font\titlefont=cmr12 at 13pt\titlefont $\bullet$ 07:58: Implement Trie}\smallskip
I forget, can the Trie be a recursive structure? It's a node with a value, a terminal node flag, and a set of children in a lookup table. I think that can be done recursively.
\par

\xrdef{node42}
{\medskip\noindent\bf\font\titlefont=cmss17 at 20pt\titlefont 322 coin change}\medskip
{\medskip\noindent\bf\font\titlefont=cmitt12 at 17pt\titlefont 2025-05-14}\medskip
{\smallskip\noindent\bf\font\titlefont=cmr12 at 13pt\titlefont $\bullet$ 08:06: coin change (initial)}\smallskip
it's too early for coin change
\par

this is a dynamic programming problem, optimizing for fewest number of coins that make up a particular amount.
\par

We are always trying to reach a particular target amount using a given set of coins. If the target amount is less than the smallest coin, there is no possible combination.
\par

When a coin does fit into the target amount, we subtract that from the target amount, and run it again, adding 1 coin. It would look something like: nCoins(amt) = 1 + nCoins(amt - coin)
\par

I have to go, but things that are tripping me up are: how to notate the recurrance relation? and how to keep track of the minimum? how to track the invalid (negative) value?
\par

{\medskip\noindent\bf\font\titlefont=cmitt12 at 17pt\titlefont 2025-05-15}\medskip
{\smallskip\noindent\bf\font\titlefont=cmr12 at 13pt\titlefont $\bullet$ 07:37: more coin change attempts}\smallskip
Working this out on pen and paper. Going to figure out the recurrance relation. If I can grok that, I feel like I'll know this problem cold.
\par

I'm looking at the answer now.
\par

Revisit coin change. Still struggling with setting up the problem and seeing the patterns.
\par

\xrdef{node43}
{\medskip\noindent\bf\font\titlefont=cmss17 at 20pt\titlefont 238 product of array except self}\medskip
{\medskip\noindent\bf\font\titlefont=cmitt12 at 17pt\titlefont 2025-05-15}\medskip
{\smallskip\noindent\bf\font\titlefont=cmr12 at 13pt\titlefont $\bullet$ 08:17: Product of array except self}\smallskip
I vaguely recall needing to work this out on paper. I think it involves keeping track of two prefix sum arrays? One from left to right, the other from right to left.
\par

The hint mentions an upper bound of 32 bits, so that could be a hint.
\par

Got it, the answer involved me carefully populating a "left" product sum array with a "right" product sum array. Writing an initial solution.
\par

\xrdef{node44}
{\medskip\noindent\bf\font\titlefont=cmss17 at 20pt\titlefont 155 min stack}\medskip
{\medskip\noindent\bf\font\titlefont=cmitt12 at 17pt\titlefont 2025-05-16}\medskip
{\smallskip\noindent\bf\font\titlefont=cmr12 at 13pt\titlefont $\bullet$ 08:06: min stack}\smallskip
I think I remember looking this one up before. I think it involves 2 stacks, one for the min value, one for the items on the stack.
\par

\xrdef{node45}
{\medskip\noindent\bf\font\titlefont=cmss17 at 20pt\titlefont 98 validate binary search tree}\medskip
{\medskip\noindent\bf\font\titlefont=cmitt12 at 17pt\titlefont 2025-05-16}\medskip
{\smallskip\noindent\bf\font\titlefont=cmr12 at 13pt\titlefont $\bullet$ 08:14: validate binary search tree}\smallskip
Looks like recursive DFS.
\par

I've hit an edge case.
\par

Right, the issue is that for a root node "5", it has a right child of 6, with a left child of 3. 3 is in the wrong place. There needs to be a minimum boundary and a maximum boundary.
\par

Looking at the answer.
\par

oh jeez, inorder traversal. I forgot that inorder traversal yields nodes in order. Basically, you keep track of the previous value. You recursively go to the leftmost node, then go back.
\par

\xrdef{node46}
{\medskip\noindent\bf\font\titlefont=cmss17 at 20pt\titlefont 200 number of islands}\medskip
{\medskip\noindent\bf\font\titlefont=cmitt12 at 17pt\titlefont 2025-05-17}\medskip
{\smallskip\noindent\bf\font\titlefont=cmr12 at 13pt\titlefont $\bullet$ 14:27: number of islands}\smallskip
This is a BFS graph traversal problem.
\par

My first approach was to use a queue to iterate through all the items in the grid. a 0 traveling to a 1 increased the numberof islands. This didn't work for the second example.
\par

I think it might need two queues. One to traverse through islands, another to actually find islands.
\par

Islands in the island queue only consume contiguous island bits. if it hits an edge... it pushes that onto the "sea" queue? When the island queue is depleted, the sea queue searches until it finds another island. When it finds one, it pushes that onto the island queue, increases the count, and breaks. Loop breaks when both queues are empty.
\par

Same exact issue. I was only going down and two the right, maybe I need to go all 4 directions. Coding that in now.
\par

I wasn't breaking out of the sea q loop early enough. I coded up an "found island" variable that gets set when an island is found. When the for loop breaks, it will also break the while loop.
\par

Trying to rework the editorial BFS solution now.
\par

The solution can be done by sweeping through the rows and columns in a for loop. If a 1 is encountered, increment the number of islands. put that on a queue and BFS the neighbors that are 1s, marking them as visited.
\par

\xrdef{node47}
{\medskip\noindent\bf\font\titlefont=cmss17 at 20pt\titlefont 994 rotting oranges}\medskip
{\medskip\noindent\bf\font\titlefont=cmitt12 at 17pt\titlefont 2025-05-17}\medskip
{\smallskip\noindent\bf\font\titlefont=cmr12 at 13pt\titlefont $\bullet$ 15:25: Rotting Oranges}\smallskip
I think I remember how to do this one. This is BFS. Initialize a queue with the first rotting oranges, their positions, and 0 days. When you iterate through the queue, traverse four directions. Keep track of the max number of days.
\par

Forgot about the impossibility scenario. Gotta count the number of oranges traversed, and compare it with the total number of oranges
\par

I was pretty sloppy this time with my code, made a lot of obvious bugs. But my aproach was pretty much sound.
\par

\xrdef{node48}
{\medskip\noindent\bf\font\titlefont=cmss17 at 20pt\titlefont 33 search in rotated sorted array}\medskip
{\medskip\noindent\bf\font\titlefont=cmitt12 at 17pt\titlefont 2025-05-20}\medskip
{\smallskip\noindent\bf\font\titlefont=cmr12 at 13pt\titlefont $\bullet$ 07:53: Search in rotated sorted array, almost got it right}\smallskip
Since it's an O(log(n)), it's a divide and conquer problem.
\par

I'd like to believe that if you segment the array, there will be always at least one that is in sorted order. We can see if the value is in range, and if it isn't, choose the other segment.
\par

lol, I got all the base cases to pass quicker than I expected. TLE for array of size 2.
\par

I'm at the point where I'm guessing and trying to tweak a solution that I believe to be mostly correct. time to look up the answer and see how far off I actually am.
\par

Damn it. The solution involved finding the pivot and then searching. Didn't do that. Will need to examine this in more detail later. But it on the TODO list with coin change.
\par

\xrdef{node49}
{\medskip\noindent\bf\font\titlefont=cmss17 at 20pt\titlefont 39 combination sum}\medskip
{\medskip\noindent\bf\font\titlefont=cmitt12 at 17pt\titlefont 2025-05-20}\medskip
{\smallskip\noindent\bf\font\titlefont=cmr12 at 13pt\titlefont $\bullet$ 08:15: combination sum}\smallskip
This feels like a kind of backtracking problem. How they use the "candidates" terminology is a hint of that.
\par

Candidates are distinct, so I don't need to deduplicate.
\par

You can use candidates any number of times. The upper limit is they have to add to target.
\par

I'm running into the "unique solution" issue. Trying to rework in such a way that each backtracking call removes candidates
\par

There needs to be a way to say okay try none of this candidate, now 1, now 2, now 3. and then recursively moving through. If that order can be achieved, then you can solve it.
\par

I am out of time.
\par

{\medskip\noindent\bf\font\titlefont=cmitt12 at 17pt\titlefont 2025-05-27}\medskip
{\smallskip\noindent\bf\font\titlefont=cmr12 at 13pt\titlefont $\bullet$ 07:56: leetcode: 39 combination sum}\smallskip
Okay, I tried to do this a few days ago. hmm. Got far enough where I was getting answers, but had issues with de-duplication.
\par

The logic I was using wasn't quite right. But I did understand that it was backtracking.
\par

\xrdef{node50}
{\medskip\noindent\bf\font\titlefont=cmss17 at 20pt\titlefont 46 permutations}\medskip
{\medskip\noindent\bf\font\titlefont=cmitt12 at 17pt\titlefont 2025-05-27}\medskip
{\smallskip\noindent\bf\font\titlefont=cmr12 at 13pt\titlefont $\bullet$ 08:21: Permutations}\smallskip
Another backtracking problem
\par

hmmm. I have the scaffolding set up. But the logic isn't quite right. Out of time.
\par

{\medskip\noindent\bf\font\titlefont=cmitt12 at 17pt\titlefont 2025-05-28}\medskip
{\smallskip\noindent\bf\font\titlefont=cmr12 at 13pt\titlefont $\bullet$ 08:08: back to permutations}\smallskip
Okay, I'm hitting the same dead end I did yesterday morning. I'm getting this problem because I'm using a "start" point, which, restricts the ordering.
\par

Doing a gross thing that I think will work, generate all the combinations, then filter the ones where the combo is unique.
\par

Gross, it works.
\par

What if we popped from nums, pushed to combok, then pushed back onto nums?
\par

Didn't work either.
\par

Okay, the solution involves backtracking, and using some kind of way to check if the number is in the combo ("if x in combo" for python). Which is annoying, because I didn't know you could do that for lists in python.
\par

\xrdef{node51}
{\medskip\noindent\bf\font\titlefont=cmss17 at 20pt\titlefont 56 merge intervals}\medskip
{\medskip\noindent\bf\font\titlefont=cmitt12 at 17pt\titlefont 2025-05-28}\medskip
{\smallskip\noindent\bf\font\titlefont=cmr12 at 13pt\titlefont $\bullet$ 08:30: merge intervals (I cheated)}\smallskip
My previous solution was still on the screen, and it was so small I saw most of it. so I'm just going to type up the solution.
\par

One thing I will say: you need to sort the intervals before processing. Python has a built-in "sorted" function, which knows to automatically sort by first value (start). That seems to do a bulk of the heavy lifting.
\par

I did try to code this out after following my own instructions, and I missed a step (append at end of loop). So, this is tricky.
\par

\xrdef{node52}
{\medskip\noindent\bf\font\titlefont=cmss17 at 20pt\titlefont 236 lowest common ancestor binary tree}\medskip
{\medskip\noindent\bf\font\titlefont=cmitt12 at 17pt\titlefont 2025-05-30}\medskip
{\smallskip\noindent\bf\font\titlefont=cmr12 at 13pt\titlefont $\bullet$ 08:08: LCA (incomplete, struggled, looked at answer)}\smallskip
I remember this one causing me issues, even after looking at the solution.
\par

The editorial I believe enumerated the possible combinations. I think there were three?
\par

It's some kind of DFS.
\par

When you do DFS, at some point you'll run into either the p or q node. One of those will come first. Somehow, you'll need to find the other node, and determine if it's a child or not.
\par

I think I'm missing some intuition here. Looking at the solution.
\par

Damn, this is the one with the recursive solution that has "left", "right", and "mid" bools. This is familiar to me, but definitely not something I would have thought up.
\par

I have to go, but next time I return to this, I will try to implement this without looking at the solution too closely. Something something recursion, return a bool. Left, right, mid. At least 2 means you've found the LCA.
\par

\xrdef{node53}
{\medskip\noindent\bf\font\titlefont=cmss17 at 20pt\titlefont news}\medskip
{\medskip\noindent\bf\font\titlefont=cmitt12 at 17pt\titlefont 2026-01-12}\medskip
{\smallskip\noindent\bf\font\titlefont=cmr12 at 13pt\titlefont $\bullet$ 17:55: digging up my leetcode notes. wonder if I can render logs from them}\smallskip
ran into rendering bugs. this may take a while. turning on the stopwatch.
\par

that's fixed and it displays. but, duplicate logs are showing up.
\par

Claude gave me a solution pretty quickly. I can use the DISTINCT keyword in the aggregate. Thanks, Claude.
\par

going to the sort the TOC alphabetically
\par

\bye
